\section{Project Structure}
\label{sec:thesis-structure}

The report is divided into seven chapters, with a brief overview of each chapter as follows:

\begin{itemize}
	\item \textbf{Chapter 1 - Introduction:} This chapter introduces the project topic, explains the motivation for a natural-language Slurm assistant, defines the scope and objectives, discusses the main challenges, and outlines the adopted solutions.

	\item \textbf{Chapter 2 - Theoretical Background:} This chapter presents the foundational knowledge required for the project, including Slurm workload management, Large Language Models, agent architectures, and the Model Context Protocol.

	\item \textbf{Chapter 3 - Related Tools and Works:} This chapter reviews existing HPC management interfaces, LLM agent frameworks, and MCP-based tool integrations, then positions the project relative to those works.

	\item \textbf{Chapter 4 - Proposed Approaches:} This chapter analyzes and designs the Slurm Agent system, including requirements, system architecture, component responsibilities, and the Observer/Operator workflow.

	\item \textbf{Chapter 5 - Implementation:} This chapter describes the implementation of the main system components: the MCP Slurm tool server, the agent backend, the frontend interface, and supporting evaluation infrastructure.

	\item \textbf{Chapter 6 - Experiments:} This chapter presents the testing approach, benchmark design, evaluation criteria, and experimental results of the Slurm Agent system.

	\item \textbf{Chapter 7 - Conclusion:} This chapter summarizes the project outcomes, discusses limitations, and reflects on the contributions relative to the project goals.
\end{itemize}

